\section{Analysis of Discrepancies}
\label{sec:analysis}

\subsection{Claim Verdicts}
\todo{Supported / partially supported / not supported / inconclusive for C1--C3.}

\subsection{Gaps}
\todo{Absolute and relative differences from the paper.}

\subsection{Evidence}
\todo{Controlled re-runs or checks supporting each explanation.}

\subsection{Limitations}
\todo{What this reproduction cannot show.}

\subsection{Race Detection on the Shared-Memory Cases}

Tu et al.~\cite{tu2019go} reported in Table~12 that Go's race
detector detected 10 of 20 non-blocking bugs, including 3 of 4
bugs involving anonymous functions. The authors counted a bug as
detected if it produced at least one report within 100 runs.

In our reproduction, the buggy Docker \#27037 variant produced
a race warning in all three runs. Each run counted once even
when the detector issued multiple warnings. This result shows
that the detector can identify the conflicting accesses to the
loop variable in our reproduction of the anonymous-function bug.
Figure~\ref{fig:race-detection} uses two units of analysis: the
paper counts distinct bugs, while our experiment counts repeated
runs of each variant. Our three runs establish repeatability
for this case under the tested setup. Estimating detection
across the paper's benchmark would require testing its broader
set of bugs under the same run protocol.

The buggy Docker \#25384 variant timed out in all three runs
without emitting a race warning. Its failure arose because
\texttt{Wait()} blocked before the loop could start the
remaining task. The test timeout exposed this blocking failure.
Table~12 covers non-blocking bugs, so this case is outside the
scope of its race-detection evaluation. Both fixed variants
passed all three tests without race warnings under the tested
workloads.

\subsection{Deviation Attribution}
\label{sec:deviation-attribution}

This subsection records differences among hosts, detectors, and language
versions. It does not replace the shared-memory race-detection discussion
above, and it does not fill the paper-versus-ours cells for claims C1--C3,
which remain unset because we have not re-extracted the paper's original
aggregate numbers from a verified source in this revision.

\paragraph{Hosts.}
The official Week~2 full-suite evidence was collected on WSL2
(Linux 6.18.58.1-microsoft-standard-WSL2, 16 available processors,
\texttt{GOMAXPROCS} unset) inside the shared Ubuntu 18.04.6 + Go 1.10.8
image, as summarised in \Cref{tab:env-official}. An earlier single-module
snapshot used Apple M2 hardware with \texttt{linux/amd64} emulation
(\Cref{tab:env-early}). The Week~3 extension, including a fresh
configuration~A baseline, ran on a native x86\_64 Debian GNU/Linux 13.7
host (kernel 6.12.113+, 8 CPUs, Docker 29.10.0) with the same Ubuntu
18.04.6 containers and no architecture emulation. On this 8-CPU host,
configuration~A reproduced the same official-checker symptoms as the WSL2
Week~2 run in all 10 repeats of all six cases. Those WSL2 counts remain
useful context, but they are not the paired A baseline for the A/B/C
comparison.

\paragraph{Timeout versus the deadlock detector.}
The two blocking cases (Kubernetes \#5316 and buggy Docker \#25384) were
exposed by \texttt{go test -timeout=5s}, not by Go's runtime deadlock
detector. None of the official Week~2 logs, and none of the 180 official
Week~3 extension logs, contains the string
\texttt{all goroutines are asleep - deadlock}. A timeout and a runtime
deadlock report are therefore not interchangeable in this reproduction.

\paragraph{Checker script BOM in the Week~2 channel logs.}
The official Week~2 channel logs still contain the line
\texttt{check\_result.sh: line 1: \#!/usr/bin/env: No such file or directory},
which is consistent with a UTF-8 byte-order mark at the start of
\texttt{bugs/channel/check\_result.sh} at the time of that run. The checker
verdicts in those logs were nonetheless correct. The current file on
\texttt{main} has a normal shebang and no BOM. The Week~3 extension invoked
that current script; none of its 180 official logs repeats the missing-interpreter
message.

\paragraph{Language version versus unverified toolchain claims.}
For Figure~8, configurations A and B both produced a race warning in 10/10
runs and never printed the unique \texttt{v1.17}--\texttt{v1.20} set, while
configuration C produced no warning in 10/10 runs and printed that unique
set every time. Because A and B already agree, and only the \texttt{go.mod}
language version changes from B to C, we attribute the disappearance of the
race to the Go 1.22 per-iteration rule. Under configuration~C the Week~2
official checker still labels all 10 buggy Figure~8 runs
\texttt{UNEXPECTED\_OR\_NOT\_REPRODUCED} because it continues to expect
\texttt{EXPECTED\_RACE}; that label is the expected scientific outcome of the
language-rule change, not evidence that ``all official checkers passed.''
Configuration~B often printed a
single repeated value (\texttt{v1.21} in nine of ten runs). That pattern is
consistent with a shared loop variable, but we do not treat it, or any
wall-clock difference, as evidence of a specific scheduler-mechanism change
in the newer runtime. Absence of a race warning is also not, by itself,
proof that a race cannot occur; the configuration-C conclusion relies on the
controlled language-version switch together with the stable unique version
set over 10 runs.

\paragraph{Limitations.}
The extension uses $N=10$ runs rather than the paper's 100-run detection
protocol, four reduced kernels rather than the paper's full bug collection,
and the default of eight logical CPUs. Wall-clock times include compilation.
\texttt{go list -m} confirmed the module language versions, but a follow-up
package-level \texttt{go list} inside the container printed VCS-stamping
errors (exit status 128) that we record without treating them as
language-version failures. Claim verdicts C1--C3 in \Cref{tab:paper-vs-ours}
remain to do.
